\section{Part 2：训练"Thinking"模型}

\begin{figure}[H]
\centering
\includegraphics[width=0.9\textwidth]{slides-images/slide-30.jpg}
\caption{Part 2：训练"Thinking"模型的主要论文\protect\footnotemark}
\end{figure}
\footnotetext{Slides 第 30 页。}

\subsection{DeepSeek-R1：通过 RL 激励推理能力}

DeepSeek-R1 是"thinking model"范式的标志性工作。其核心发现是：\textbf{仅通过 RL 训练（不需要人工标注的推理链），模型就能自发涌现出复杂的推理行为}。

\begin{importantbox}{DeepSeek-R1 的训练流程}
\begin{enumerate}
    \item 在基础模型上直接应用 GRPO（Group Relative Policy Optimization）进行 RL 训练
    \item 奖励仅基于最终答案是否正确（outcome reward）+ 格式奖励
    \item 模型自发学会了：验证（verification）、回溯（backtracking）、设定子目标（subgoal setting）、反向推理（backward chaining）等元策略
\end{enumerate}
\end{importantbox}

\begin{knowledgebox}{GRPO（Group Relative Policy Optimization）}
GRPO 是 PPO 的简化变体。对于每个问题采样一组解答，用组内的相对奖励（减去组内均值后标准化）作为优势估计，避免了训练额外的 critic/value 网络。这在大规模 LLM 训练中更加高效。
\end{knowledgebox}

\subsection{从具体轨迹理解 ``thinking'' 是什么}

``Thinking model'' 这个说法很容易被说得神秘，但讲者展示的例子其实很具体：模型会在一次回答内部执行更多宏观操作，例如尝试一个方向、发现不对、显式验证、回退并重规划。也就是说，RL 并不是单纯让输出更长，而是提高了模型在固定预算内组织这些宏观动作的能力。

\begin{figure}[H]
\centering
\includegraphics[width=0.9\textwidth]{slides-images/slide-31.jpg}
\caption{Thinking model 的单条轨迹示意：验证、回溯与重新规划交替出现\protect\footnotemark}
\end{figure}
\footnotetext{Slides 第 31 页。}

\subsection{真正改变的是动作空间，而非 RL 教科书}

讲者在 slides 中专门给出一句很重要的判断：\textbf{训练目标并没有本质改变，真正变化的是模型已经能执行的 ``macro actions''。} DeepSeek-R1 仍在做 policy gradient，Kimi K1.5 也仍然依赖 advantage 驱动的策略更新，但基础模型已经能以更长 token budget 表达验证、回溯、子目标分解等高阶行为，因此同样的 RL 配方能在新的动作空间里产生更强效果。

\begin{figure}[H]
\centering
\includegraphics[width=0.9\textwidth]{slides-images/slide-32.jpg}
\caption{Thinking model 时代的变化：奖励形式近似不变，动作空间和预算显著扩大\protect\footnotemark}
\end{figure}
\footnotetext{Slides 第 32 页。}

\begin{knowledgebox}{为什么 ``动作空间升级'' 比 ``目标函数升级'' 更关键}
很多讨论把 thinking model 的突破归因于某个新优化器或新损失，但这节课的判断更克制：只要基础模型已经具备一定的元认知雏形，扩大上下文预算、允许更长 rollout、再用 RL 奖励这些行为，就足以显著放大推理表现。换句话说，\textbf{新能力往往来自已有能力的放大和重组，而不是凭空发明一种全新训练目标。}
\end{knowledgebox}

\subsection{涌现的推理行为}

\begin{figure}[H]
\centering
\includegraphics[width=0.9\textwidth]{slides-images/slide-35.jpg}
\caption{``Action'' Space：涌现的元策略对比\protect\footnotemark}
\end{figure}
\footnotetext{Slides 第 35 页。}

Gandhi et al. (2025) 的研究表明，能否在 RL 训练中涌现认知行为取决于基础模型是否已经具备这些行为的"种子"。DeepSeek 系列模型在 RL 训练过程中，验证和回溯行为的频率先上升后下降；而 Meta 的 Llama 模型则完全没有展现出这些行为，RL 训练也未能改善其推理能力。

\begin{warningbox}{RL 并非万能}
RL 无法让模型学会它"完全不会"的行为。如果基础模型在预训练阶段没有接触过推理链式思考的数据，单靠 RL 很难从零涌现出复杂推理行为。基础模型的能力是 RL 发挥作用的前提。
\end{warningbox}

\subsection{固定预算下的训练目标}

讲者随后把问题重新表述为一个很工程化的目标：在每道测试题拥有固定总 compute budget 的前提下，训练一个策略，让 sampled response 在这个预算内尽可能提高成功率。这个表述有两个现实含义：

\begin{enumerate}
    \item 推理不是 ``越长越好''，而是要在\textbf{受限预算}内决定把算力花在哪里。
    \item RL、SFT、RFT 本质上都可以被看作这个适应问题的不同求解器，它们差异主要在是否利用奖励以及利用奖励的方式有多充分。
\end{enumerate}

\begin{figure}[H]
\centering
\includegraphics[width=0.9\textwidth]{slides-images/slide-33.jpg}
\caption{Thinking model 训练的目标：在固定测试预算内最大化成功率\protect\footnotemark}
\end{figure}
\footnotetext{Slides 第 33 页。}

\begin{figure}[H]
\centering
\includegraphics[width=0.9\textwidth]{slides-images/slide-34.jpg}
\caption{同一预算约束下，RL 与 SFT/RFT 都可视为 ``适应问题'' 的不同实现\protect\footnotemark}
\end{figure}
\footnotetext{Slides 第 34 页。}

\subsection{Kimi K1.5：扩展 RL 训练}

Kimi K1.5 的核心贡献在于如何在大规模上有效地进行 RL 训练：
\begin{itemize}
    \item 更长的上下文窗口（128K tokens）
    \item 更好的采样策略和奖励设计
    \item 实验表明 RL 训练的性能随计算量增加而持续提升
\end{itemize}

\subsection{长度放大与 test-time scaling 的新范式}

RL 训练的另一个直接后果，是模型愿意把更多 token 预算用在真正有价值的地方。讲者引用 Kimi 的报告指出，RL 训练后模型回答长度显著增长，但这不是无意义冗长，而是伴随着更多检查、修正和分支探索。于是，``推理时多给一些计算'' 开始成为一条可与 ``继续增大模型参数'' 并列的性能提升路线。

\begin{figure}[H]
\centering
\includegraphics[width=0.9\textwidth]{slides-images/slide-36.jpg}
\caption{RL 训练会放大回答长度，并把额外长度转化为更丰富的推理行为\protect\footnotemark}
\end{figure}
\footnotetext{Slides 第 36 页。}

\begin{figure}[H]
\centering
\includegraphics[width=0.9\textwidth]{slides-images/slide-37.jpg}
\caption{新的问题设定：除了放大模型，也可以放大 test-time compute\protect\footnotemark}
\end{figure}
\footnotetext{Slides 第 37 页。}

\begin{figure}[H]
\centering
\includegraphics[width=0.9\textwidth]{slides-images/slide-39.jpg}
\caption{更长的思考、更强的自校正构成了新的 test-time scaling 路线\protect\footnotemark}
\end{figure}
\footnotetext{Slides 第 39 页。}

\subsection{Thinking 模型的效果}

\begin{figure}[H]
\centering
\includegraphics[width=0.9\textwidth]{slides-images/slide-40.jpg}
\caption{Thinking 模型在各 benchmark 上的表现\protect\footnotemark}
\end{figure}
\footnotetext{Slides 第 40 页。}

Thinking 模型在 Codeforces 竞赛编程、Humanity's Last Exam、AIME 等多个基准上取得了显著提升。DeepSeek-R1 在 AIME 2024 上达到 79.8\% 的 Pass@1，与 OpenAI o1-1217 的 79.2\% 相当。

\subsection{Test-Time Compute 与 Meta-RL}

讲者最后提到，优化 test-time compute（推理时计算量分配）本质上是一个 Meta-RL 问题：模型需要学习在推理时如何"花费"计算预算。

\subsection{还没有解决的问题}

尽管 leaderboard 成绩已经很亮眼，这节课最后一页反而是警告意味最强的一页：当前 thinking model 仍然缺乏统一而清晰的 formulation。我们还没有彻底回答三件事：第一，是否存在比 outcome reward 更稳健的 dense reward 设计；第二，在什么条件下 RL 的收益会稳定超过 SFT；第三，如何把 ``推理增强'' 严格表述为一个适应问题，而不是一组经验技巧的堆叠。

\begin{figure}[H]
\centering
\includegraphics[width=0.9\textwidth]{slides-images/slide-41.jpg}
\caption{课程结尾给出的开放问题：formulation、dense rewards 与 RL 对 SFT 的优势边界\protect\footnotemark}
\end{figure}
\footnotetext{Slides 第 41 页。}

\begin{warningbox}{真正的挑战不在于复刻一篇 DeepSeek-R1}
如果只盯着某个公开模型或某条 benchmark 曲线，很容易把问题误解成 ``找到同款 recipe''。但 Aviral Kumar 的收束方式更像研究问题清单：我们仍然需要弄清楚奖励信号的形状、搜索与验证的耦合方式、以及不同基础模型为何对 RL 的响应差异如此巨大。
\end{warningbox}

\subsection{本章小结}

DeepSeek-R1 等 thinking model 证明了：(1) RL 可以激励 LLM 自主发展推理能力；(2) 关键前提是基础模型需具备推理行为的"种子"；(3) GRPO 等简化的策略优化方法在大规模训练中既高效又有效。

